% Options for packages loaded elsewhere
\PassOptionsToPackage{unicode}{hyperref}
\PassOptionsToPackage{hyphens}{url}
%
\documentclass[
]{article}
\usepackage{amsmath,amssymb}
\usepackage{iftex}
\ifPDFTeX
  \usepackage[T1]{fontenc}
  \usepackage[utf8]{inputenc}
  \usepackage{textcomp} % provide euro and other symbols
\else % if luatex or xetex
  \usepackage{unicode-math} % this also loads fontspec
  \defaultfontfeatures{Scale=MatchLowercase}
  \defaultfontfeatures[\rmfamily]{Ligatures=TeX,Scale=1}
\fi
\usepackage{lmodern}
\ifPDFTeX\else
  % xetex/luatex font selection
\fi
% Use upquote if available, for straight quotes in verbatim environments
\IfFileExists{upquote.sty}{\usepackage{upquote}}{}
\IfFileExists{microtype.sty}{% use microtype if available
  \usepackage[]{microtype}
  \UseMicrotypeSet[protrusion]{basicmath} % disable protrusion for tt fonts
}{}
\makeatletter
\@ifundefined{KOMAClassName}{% if non-KOMA class
  \IfFileExists{parskip.sty}{%
    \usepackage{parskip}
  }{% else
    \setlength{\parindent}{0pt}
    \setlength{\parskip}{6pt plus 2pt minus 1pt}}
}{% if KOMA class
  \KOMAoptions{parskip=half}}
\makeatother
\usepackage{xcolor}
\usepackage{color}
\usepackage{fancyvrb}
\newcommand{\VerbBar}{|}
\newcommand{\VERB}{\Verb[commandchars=\\\{\}]}
\DefineVerbatimEnvironment{Highlighting}{Verbatim}{commandchars=\\\{\}}
% Add ',fontsize=\small' for more characters per line
\newenvironment{Shaded}{}{}
\newcommand{\AlertTok}[1]{\textcolor[rgb]{1.00,0.00,0.00}{\textbf{#1}}}
\newcommand{\AnnotationTok}[1]{\textcolor[rgb]{0.38,0.63,0.69}{\textbf{\textit{#1}}}}
\newcommand{\AttributeTok}[1]{\textcolor[rgb]{0.49,0.56,0.16}{#1}}
\newcommand{\BaseNTok}[1]{\textcolor[rgb]{0.25,0.63,0.44}{#1}}
\newcommand{\BuiltInTok}[1]{\textcolor[rgb]{0.00,0.50,0.00}{#1}}
\newcommand{\CharTok}[1]{\textcolor[rgb]{0.25,0.44,0.63}{#1}}
\newcommand{\CommentTok}[1]{\textcolor[rgb]{0.38,0.63,0.69}{\textit{#1}}}
\newcommand{\CommentVarTok}[1]{\textcolor[rgb]{0.38,0.63,0.69}{\textbf{\textit{#1}}}}
\newcommand{\ConstantTok}[1]{\textcolor[rgb]{0.53,0.00,0.00}{#1}}
\newcommand{\ControlFlowTok}[1]{\textcolor[rgb]{0.00,0.44,0.13}{\textbf{#1}}}
\newcommand{\DataTypeTok}[1]{\textcolor[rgb]{0.56,0.13,0.00}{#1}}
\newcommand{\DecValTok}[1]{\textcolor[rgb]{0.25,0.63,0.44}{#1}}
\newcommand{\DocumentationTok}[1]{\textcolor[rgb]{0.73,0.13,0.13}{\textit{#1}}}
\newcommand{\ErrorTok}[1]{\textcolor[rgb]{1.00,0.00,0.00}{\textbf{#1}}}
\newcommand{\ExtensionTok}[1]{#1}
\newcommand{\FloatTok}[1]{\textcolor[rgb]{0.25,0.63,0.44}{#1}}
\newcommand{\FunctionTok}[1]{\textcolor[rgb]{0.02,0.16,0.49}{#1}}
\newcommand{\ImportTok}[1]{\textcolor[rgb]{0.00,0.50,0.00}{\textbf{#1}}}
\newcommand{\InformationTok}[1]{\textcolor[rgb]{0.38,0.63,0.69}{\textbf{\textit{#1}}}}
\newcommand{\KeywordTok}[1]{\textcolor[rgb]{0.00,0.44,0.13}{\textbf{#1}}}
\newcommand{\NormalTok}[1]{#1}
\newcommand{\OperatorTok}[1]{\textcolor[rgb]{0.40,0.40,0.40}{#1}}
\newcommand{\OtherTok}[1]{\textcolor[rgb]{0.00,0.44,0.13}{#1}}
\newcommand{\PreprocessorTok}[1]{\textcolor[rgb]{0.74,0.48,0.00}{#1}}
\newcommand{\RegionMarkerTok}[1]{#1}
\newcommand{\SpecialCharTok}[1]{\textcolor[rgb]{0.25,0.44,0.63}{#1}}
\newcommand{\SpecialStringTok}[1]{\textcolor[rgb]{0.73,0.40,0.53}{#1}}
\newcommand{\StringTok}[1]{\textcolor[rgb]{0.25,0.44,0.63}{#1}}
\newcommand{\VariableTok}[1]{\textcolor[rgb]{0.10,0.09,0.49}{#1}}
\newcommand{\VerbatimStringTok}[1]{\textcolor[rgb]{0.25,0.44,0.63}{#1}}
\newcommand{\WarningTok}[1]{\textcolor[rgb]{0.38,0.63,0.69}{\textbf{\textit{#1}}}}
\usepackage{longtable,booktabs,array}
\usepackage{calc} % for calculating minipage widths
% Correct order of tables after \paragraph or \subparagraph
\usepackage{etoolbox}
\makeatletter
\patchcmd\longtable{\par}{\if@noskipsec\mbox{}\fi\par}{}{}
\makeatother
% Allow footnotes in longtable head/foot
\IfFileExists{footnotehyper.sty}{\usepackage{footnotehyper}}{\usepackage{footnote}}
\makesavenoteenv{longtable}
\setlength{\emergencystretch}{3em} % prevent overfull lines
\providecommand{\tightlist}{%
  \setlength{\itemsep}{0pt}\setlength{\parskip}{0pt}}
\setcounter{secnumdepth}{-\maxdimen} % remove section numbering
\ifLuaTeX
  \usepackage{selnolig}  % disable illegal ligatures
\fi
\usepackage{bookmark}
\IfFileExists{xurl.sty}{\usepackage{xurl}}{} % add URL line breaks if available
\urlstyle{same}
\hypersetup{
  pdftitle={Active Cost-Sensitive Monitoring of Black-Box LLM Behavior},
  hidelinks,
  pdfcreator={LaTeX via pandoc}}

\title{Active Cost-Sensitive Monitoring of Black-Box LLM Behavior}
\author{}
\date{}

\begin{document}
\maketitle

\section{Active Cost-Sensitive Monitoring of Black-Box LLM
Behavior}\label{active-cost-sensitive-monitoring-of-black-box-llm-behavior}

\subsection{4. Probabilistic View}\label{probabilistic-view}

\subsubsection{4.1 Problem as a partially observed decision
process}\label{problem-as-a-partially-observed-decision-process}

A hosted language-model API exposes outputs and some runtime
measurements but does not expose the provider's internal model state,
deployment history, or the exact cause of a behavioral change. The
monitoring problem can therefore be represented as a partially observed
decision problem. Let the hidden state be

{[} S\_t
\in \{\text{STABLE},\text{DEGRADED},\text{TRANSIENT},\text{DISTRIBUTION\_SHIFT}\}.
{]}

The monitor receives an observation

{[} O\_t=(q,d,f,s,l,e,r,c), {]}

where \texttt{q} is task quality, \texttt{d} is semantic drift,
\texttt{f} is formatting failure, \texttt{s} is a safety-shift
indicator, \texttt{l} is latency, \texttt{e} is API error rate,
\texttt{r} is retained feedback rate, and \texttt{c} is an optional
capability-specific probe result. \texttt{feedback\_rate} exists in the
implementation as a retained field but is not currently used in belief
updates. \texttt{capability\_issue} is observable only through the
targeted capability probe.

The agent must choose one of three actions:

\begin{longtable}[]{@{}
  >{\raggedright\arraybackslash}p{(\columnwidth - 2\tabcolsep) * \real{0.5000}}
  >{\raggedright\arraybackslash}p{(\columnwidth - 2\tabcolsep) * \real{0.5000}}@{}}
\toprule\noalign{}
\begin{minipage}[b]{\linewidth}\raggedright
Action
\end{minipage} & \begin{minipage}[b]{\linewidth}\raggedright
Meaning
\end{minipage} \\
\midrule\noalign{}
\endhead
\bottomrule\noalign{}
\endlastfoot
\texttt{ACCEPT} & Continue normal operation. \\
\texttt{INVESTIGATE} & Spend an additional evidence budget before taking
the final action. \\
\texttt{REJECT} & Block, quarantine, or otherwise stop the current
behavior. \\
\end{longtable}

The hidden state is used by the simulator for evaluation, not by the
agent at decision time.

\subsubsection{4.2 Hidden state semantics}\label{hidden-state-semantics}

The four-state model intentionally separates persistent degradation,
short-lived anomalies, and input-distribution changes. The states are
not intended to identify the provider's internal cause with certainty.

\begin{longtable}[]{@{}
  >{\raggedright\arraybackslash}p{(\columnwidth - 2\tabcolsep) * \real{0.5000}}
  >{\raggedright\arraybackslash}p{(\columnwidth - 2\tabcolsep) * \real{0.5000}}@{}}
\toprule\noalign{}
\begin{minipage}[b]{\linewidth}\raggedright
State
\end{minipage} & \begin{minipage}[b]{\linewidth}\raggedright
Operational interpretation
\end{minipage} \\
\midrule\noalign{}
\endhead
\bottomrule\noalign{}
\endlastfoot
\texttt{STABLE} & Behavior remains consistent with the trusted reference
condition. \\
\texttt{DEGRADED} & Persistent loss of expected capability. \\
\texttt{TRANSIENT} & Temporary anomaly or short-lived disruption. \\
\texttt{DISTRIBUTION\_SHIFT} & The request population changes in a way
that affects observed behavior. \\
\end{longtable}

The simulator explicitly generates all four states. The limitation is
not state generation but empirical validation: the probabilities and
evidence likelihoods are synthetic assumptions and have not been
calibrated against production traces.

\subsubsection{4.3 Prior distribution}\label{prior-distribution}

The initial belief is deliberately hypothetical and is the same for
every case:

\begin{longtable}[]{@{}lr@{}}
\toprule\noalign{}
Hidden state & Prior probability \\
\midrule\noalign{}
\endhead
\bottomrule\noalign{}
\endlastfoot
\texttt{STABLE} & 0.72 \\
\texttt{DEGRADED} & 0.12 \\
\texttt{TRANSIENT} & 0.08 \\
\texttt{DISTRIBUTION\_SHIFT} & 0.08 \\
\textbf{Total} & \textbf{1.00} \\
\end{longtable}

The prior is a reproducibility device, not an empirical estimate. It
should not be interpreted as the expected frequency of these states in a
real deployment.

\subsubsection{4.4 Evidence model}\label{evidence-model}

The current model uses a naive-Bayes-style update. For evidence items
(E\_1,\ldots,E\_k), the posterior is proportional to

{[} P(S=s\mid E\_1,\ldots,E\_k)
\propto P(S=s)\prod\_\{i=1\}\^{}\{k\}P(E\_i\mid S=s). {]}

This assumes that the evidence streams are conditionally independent
given the hidden state. That assumption is convenient and auditable, but
it is not demonstrated. For example, low quality, semantic drift,
formatting problems, and API errors can be correlated during a single
failure event.

The synthetic likelihood table used by the implementation is:

\begin{longtable}[]{@{}lrrrr@{}}
\toprule\noalign{}
Evidence event & STABLE & DEGRADED & TRANSIENT & DISTRIBUTION\_SHIFT \\
\midrule\noalign{}
\endhead
\bottomrule\noalign{}
\endlastfoot
Low output quality & 0.08 & 0.78 & 0.60 & 0.35 \\
Semantic drift & 0.10 & 0.72 & 0.58 & 0.70 \\
Formatting failure & 0.04 & 0.55 & 0.35 & 0.22 \\
Safety shift & 0.03 & 0.38 & 0.18 & 0.30 \\
High latency & 0.10 & 0.33 & 0.72 & 0.15 \\
High API error & 0.04 & 0.45 & 0.55 & 0.12 \\
Capability issue probe & 0.01 & 0.95 & 0.10 & 0.20 \\
\end{longtable}

All seven columns in this table are synthetic hypotheses. They are
inspectable in the implementation and are not presented as calibrated
production probabilities.

\subsubsection{4.5 Worked posterior
update}\label{worked-posterior-update}

A worked example uses the following observed evidence:

\begin{itemize}
\tightlist
\item
  low output quality
\item
  semantic drift
\item
  no formatting failure
\item
  no safety shift
\item
  normal latency
\item
  no API error
\end{itemize}

For \texttt{DEGRADED}, the unnormalized contribution is calculated as:

\begin{Shaded}
\begin{Highlighting}[]
\NormalTok{0.12 x 0.78 = 0.093600}
\NormalTok{0.093600 x 0.72 = 0.067392}
\NormalTok{0.067392 x 0.45 = 0.030326}
\NormalTok{0.030326 x 0.62 = 0.018802}
\NormalTok{0.018802 x 0.67 = 0.012598}
\NormalTok{0.012598 x 0.55 = 0.006929}
\end{Highlighting}
\end{Shaded}

The full set of state contributions is:

\begin{longtable}[]{@{}lrr@{}}
\toprule\noalign{}
State & Prior & Joint contribution \\
\midrule\noalign{}
\endhead
\bottomrule\noalign{}
\endlastfoot
\texttt{STABLE} & 0.72 & 0.004634 \\
\texttt{DEGRADED} & 0.12 & 0.006929 \\
\texttt{TRANSIENT} & 0.08 & 0.001870 \\
\texttt{DISTRIBUTION\_SHIFT} & 0.08 & 0.008005 \\
\textbf{Evidence probability} & & \textbf{0.021437} \\
\end{longtable}

Normalizing by 0.021437 gives:

\begin{longtable}[]{@{}lr@{}}
\toprule\noalign{}
State & Posterior probability \\
\midrule\noalign{}
\endhead
\bottomrule\noalign{}
\endlastfoot
\texttt{STABLE} & 21.6\% \\
\texttt{DEGRADED} & 32.3\% \\
\texttt{TRANSIENT} & 8.7\% \\
\texttt{DISTRIBUTION\_SHIFT} & 37.3\% \\
\textbf{Total} & \textbf{100.0\%} \\
\end{longtable}

The evidence therefore changes the most likely explanation from
\texttt{STABLE} under the prior to \texttt{DISTRIBUTION\_SHIFT} under
the posterior, but it does not produce a concentrated belief. Several
states remain plausible.

\subsection{5. Information Theory}\label{information-theory}

\subsubsection{5.1 Entropy is not the same as
confidence}\label{entropy-is-not-the-same-as-confidence}

Belief uncertainty is measured with Shannon entropy:

{[} H(B)=-\sum\_s P(s)\log\_2 P(s). {]}

For the prior ((0.72,0.12,0.08,0.08)),

{[} H\_\{before\}=1.291\text{ bits}. {]}

For the worked posterior ((0.216,0.323,0.087,0.373)),

{[} H\_\{after\}=1.842\text{ bits}. {]}

Therefore,

{[} \Delta H=H\_\{after\}-H\_\{before\}=1.842-1.291=+0.551\text{ bits}.
{]}

The uncertainty increased. This is not an arithmetic error. The evidence
weakened the dominant \texttt{STABLE} hypothesis without strongly
establishing a single alternative, so the belief mass became more
distributed across multiple states.

For a realized evidence bundle, a negative information gain is possible
because the particular observation can be more ambiguous than the prior.
The active probe-selection problem is different: it uses
\textbf{expected} information gain before observing the probe outcome.

\subsubsection{5.2 Expected information gain for candidate
probes}\label{expected-information-gain-for-candidate-probes}

For a candidate binary probe (X), expected information gain is

{[} EIG(X)=H(B)-\left[P(X=1)H(B\mid X=1)+P(X=0)H(B\mid X=0)\right]. {]}

Using the synthetic prior and the likelihood table, the candidate
signals have the following expected information gains:

\begin{longtable}[]{@{}
  >{\raggedright\arraybackslash}p{(\columnwidth - 6\tabcolsep) * \real{0.2000}}
  >{\raggedleft\arraybackslash}p{(\columnwidth - 6\tabcolsep) * \real{0.2667}}
  >{\raggedleft\arraybackslash}p{(\columnwidth - 6\tabcolsep) * \real{0.2667}}
  >{\raggedleft\arraybackslash}p{(\columnwidth - 6\tabcolsep) * \real{0.2667}}@{}}
\toprule\noalign{}
\begin{minipage}[b]{\linewidth}\raggedright
Candidate probe
\end{minipage} & \begin{minipage}[b]{\linewidth}\raggedleft
P(probe positive)
\end{minipage} & \begin{minipage}[b]{\linewidth}\raggedleft
Expected information gain, bits
\end{minipage} & \begin{minipage}[b]{\linewidth}\raggedleft
EIG / 0.8 cost units
\end{minipage} \\
\midrule\noalign{}
\endhead
\bottomrule\noalign{}
\endlastfoot
\texttt{capability\_issue} & 0.1452 & 0.4099 & 0.5124 \\
\texttt{quality\_low} & 0.2272 & 0.2399 & 0.2999 \\
\texttt{semantic\_drift} & 0.2608 & 0.2386 & 0.2983 \\
\texttt{error\_high} & 0.1364 & 0.1594 & 0.1992 \\
\texttt{format\_failure} & 0.1404 & 0.1562 & 0.1952 \\
\texttt{latency\_high} & 0.1812 & 0.1180 & 0.1475 \\
\texttt{safety\_shift} & 0.1056 & 0.1067 & 0.1333 \\
\end{longtable}

The division by 0.8 is only an illustrative
information-per-investigation-unit ratio because the project does not
yet measure source-specific monetary cost or probe time. It is not a
claim that one bit has a stable monetary value.

Two design details matter here. First, \texttt{quality\_low} is already
observed in the initial observation, so it is not a discretionary probe
choice. Second, \texttt{capability\_issue} has the highest expected
information gain in this synthetic prior but is reserved for the
dedicated capability-coverage audit rather than the discretionary
candidate set. The active discretionary selector currently considers
semantic drift, formatting failure, safety shift, latency, and API
error.

\subsubsection{5.3 Value of information and entropy are
complementary}\label{value-of-information-and-entropy-are-complementary}

The project does not equate information gain with decision value. A
probe can reduce entropy without changing the preferred final action.
Conversely, a probe can have modest entropy reduction but still be
valuable if it often moves the posterior across the cost-sensitive
accept/reject boundary.

The policy therefore uses one-step \textbf{decision-theoretic value of
information} rather than entropy reduction alone:

{[} VOI\_\{gross\}=C\_\{direct\}-E{[}C\_\{post\mbox{-}probe\}{]}. {]}

A probe is worth purchasing only if the gross reduction in expected
direct decision loss is greater than its investigation cost.

\subsection{6. Evidence Selection}\label{evidence-selection}

\subsubsection{6.1 Candidate set}\label{candidate-set}

The discretionary probe selector considers five evidence sources:

\begin{enumerate}
\def\labelenumi{\arabic{enumi}.}
\tightlist
\item
  semantic drift
\item
  formatting failure
\item
  safety shift
\item
  high latency
\item
  high API error
\end{enumerate}

For each candidate, the model computes the expected direct decision loss
after a positive probe result and after a negative probe result, weights
the two cases by their predicted probabilities, and compares the
resulting expected loss with the current direct loss.

The selector does not use a generic rule such as ``choose the
highest-entropy signal.'' It uses the signal's expected impact on the
final decision.

\subsubsection{6.2 Why a capability probe was added
separately}\label{why-a-capability-probe-was-added-separately}

Failure analysis exposed a structural problem: a narrow degraded
capability could coexist with healthy average quality, low semantic
drift, normal latency, and low API errors. In such a case, no threshold
adjustment can recover information that is not present in the
observation layer.

The concrete design change was therefore a bounded capability-coverage
audit. Every fourth observation within a task category triggers a
\texttt{capability\_issue} probe, subject to a maximum of two
investigation rounds.

The capability probe has the following synthetic likelihoods:

\begin{longtable}[]{@{}lr@{}}
\toprule\noalign{}
Hidden state & P(capability\_issue) \\
\midrule\noalign{}
\endhead
\bottomrule\noalign{}
\endlastfoot
\texttt{STABLE} & 0.01 \\
\texttt{DEGRADED} & 0.95 \\
\texttt{TRANSIENT} & 0.10 \\
\texttt{DISTRIBUTION\_SHIFT} & 0.20 \\
\end{longtable}

This schedule is not claimed to be optimal. It is a targeted
intervention motivated by an observed failure class.

\subsubsection{6.3 Probe-cost knowledge
gap}\label{probe-cost-knowledge-gap}

The project currently has one common investigation cost, 0.8 decision
units. It does \textbf{not} contain empirically measured probe latency,
compute usage, API pricing, or human-review time for each evidence
source. Therefore a source-specific ``bits per rupee'' analysis would be
fabricated if written as a measured result.

The paper instead reports expected information gain and clearly labels
the 0.8 cost ratio as a synthetic sensitivity quantity. A production
version should replace it with measured cost and latency for each probe
type.

\subsection{7. Decision Policy}\label{decision-policy}

\subsubsection{7.1 Direct accept/reject
threshold}\label{direct-acceptreject-threshold}

The hypothetical action-cost model is:

\begin{longtable}[]{@{}lr@{}}
\toprule\noalign{}
Outcome & Cost \\
\midrule\noalign{}
\endhead
\bottomrule\noalign{}
\endlastfoot
Correct \texttt{ACCEPT} & 0.0 \\
Correct \texttt{REJECT} & 0.0 \\
Investigation & 0.8 \\
False \texttt{REJECT} & 3.0 \\
False \texttt{ACCEPT} & 10.0 \\
\end{longtable}

Let

{[} p=P(DEGRADED\mid E). {]}

The expected cost of accepting is

{[} C\_\{ACCEPT\}=10p. {]}

The expected cost of rejecting is

{[} C\_\{REJECT\}=3(1-p). {]}

At the break-even point,

{[} 10p=3(1-p), {]}

so

{[} 13p=3, {]}

and therefore

{[} p=\frac{3}{13}=0.230769. {]}

The implemented rule is therefore:

\begin{longtable}[]{@{}ll@{}}
\toprule\noalign{}
Posterior P(DEGRADED) & Final direct action \\
\midrule\noalign{}
\endhead
\bottomrule\noalign{}
\endlastfoot
(p\le 0.230769) & \texttt{ACCEPT} \\
(p\textgreater0.230769) & \texttt{REJECT} \\
\end{longtable}

At exactly 0.230769, the implementation assigns the tie to
\texttt{ACCEPT} because
\texttt{accept\_cost\ \textless{}=\ reject\_cost}.

This threshold is derived from the declared costs. It is not an
arbitrary 0.5 probability threshold.

\subsubsection{7.2 Correct investigation
threshold}\label{correct-investigation-threshold}

The investigation rule previously mixed two quantities and reported an
inconsistent expected cost. The revised implementation separates them:

{[} VOI\_\{gross\}=C\_\{direct\}-E{[}C\_\{post\mbox{-}probe\}{]}, {]}

{[} VOI\_\{net\}=VOI\_\{gross\}-C\_\{investigate\}. {]}

A discretionary investigation is selected only when

{[} VOI\_\{gross\}\textgreater C\_\{investigate\}. {]}

With the current hypothetical investigation cost of 0.8, a probe must
save more than 0.8 expected decision-cost units before it is purchased.

The expected cost stored for an investigation is

{[} E{[}C\_\{investigate\}{]}=E{[}C\_\{post\mbox{-}probe\}{]}+0.8. {]}

The implementation test verifies this accounting directly.

\subsubsection{7.3 Bounded investigation
budget}\label{bounded-investigation-budget}

The monitor allows at most two probe rounds. The budget exists because
an active monitor should not continue collecting evidence indefinitely
while delaying an action.

The stopping rule is:

\begin{enumerate}
\def\labelenumi{\arabic{enumi}.}
\tightlist
\item
  stop and act when gross VOI is no longer greater than the
  investigation cost, or
\item
  stop when two investigation rounds have been consumed.
\end{enumerate}

After the budget is exhausted, the same cost-derived direct
\texttt{ACCEPT}/\texttt{REJECT} rule is applied.

\subsubsection{7.4 Scheduled audits versus discretionary
VOI}\label{scheduled-audits-versus-discretionary-voi}

The coverage audit is deliberately separate from discretionary VOI. A
scheduled audit can fire because its task-category coverage condition is
due even when the current one-step VOI calculation would not purchase a
discretionary probe.

Therefore the paper makes no claim that \textbf{every} investigation is
individually justified by the VOI threshold. The precise claim is
narrower: discretionary investigations are cost-gated by VOI, while
scheduled capability audits are a bounded coverage intervention
introduced to address a demonstrated blind spot.

\subsubsection{7.5 Human-control boundary}\label{human-control-boundary}

The current implementation does not emit a \texttt{HUMAN\_REVIEW}
action. Human review is treated as an operational boundary for cases
where a real deployment would not accept an autonomous high-impact
action on the available evidence.

The prototype therefore remains a decision gate rather than an
autonomous incident commander.

\subsection{8. Experiment}\label{experiment}

\subsubsection{8.1 Primary 50-case
benchmark}\label{primary-50-case-benchmark}

The primary evaluation contains 50 deterministic synthetic cases
generated with random seed 7. Five task categories are used:

\begin{itemize}
\tightlist
\item
  \texttt{math}
\item
  \texttt{code}
\item
  \texttt{instruction}
\item
  \texttt{safety}
\item
  \texttt{knowledge}
\end{itemize}

The state schedule is deterministic. Cases 16 through 26 are
\texttt{DEGRADED}; cases 31 through 35 plus cases 43 and 44 are
\texttt{TRANSIENT}; cases 37 through 41 are
\texttt{DISTRIBUTION\_SHIFT}; all remaining cases are \texttt{STABLE}.

Two adversarial conditions are injected:

\begin{enumerate}
\def\labelenumi{\arabic{enumi}.}
\tightlist
\item
  Cases 16 through 20 are a narrow degraded capability blind spot. Their
  aggregate quality remains around 0.84, semantic drift remains low,
  formatting remains valid, latency remains near normal, and API error
  rate remains low.
\item
  Case 30 is a stable measurement-noise case whose observed quality,
  semantic drift, formatting, latency, and error rate resemble
  degradation.
\end{enumerate}

The same 50-case dataset is evaluated by three policies:

\begin{longtable}[]{@{}
  >{\raggedright\arraybackslash}p{(\columnwidth - 2\tabcolsep) * \real{0.5000}}
  >{\raggedright\arraybackslash}p{(\columnwidth - 2\tabcolsep) * \real{0.5000}}@{}}
\toprule\noalign{}
\begin{minipage}[b]{\linewidth}\raggedright
Policy
\end{minipage} & \begin{minipage}[b]{\linewidth}\raggedright
Definition
\end{minipage} \\
\midrule\noalign{}
\endhead
\bottomrule\noalign{}
\endlastfoot
Baseline & Accept when quality is at least 0.70, otherwise reject. \\
Binary & Use the belief update and the cost-derived direct accept/reject
rule, but no active evidence collection. \\
Active & Use the belief update, corrected VOI-based investigation, a
two-round budget, and the every-fourth-observation capability coverage
audit. \\
\end{longtable}

The hidden simulator state is used only after each decision to calculate
the evaluation metrics.

\subsubsection{8.2 120-case re-test}\label{case-re-test}

A second run contains 120 cases using the same random seed, 7, and the
same 50-case state schedule repeated over the 120-case sequence. New
observations are sampled within the repeated state schedule. This is a
scaling and before/after re-test, not a statistically independent
benchmark.

The retained pre-change 120-case run had an active policy that selected
no investigations. The post-change run uses the corrected investigation
accounting and the capability-coverage audit.

\subsubsection{8.3 Evaluation metrics}\label{evaluation-metrics}

For each policy, the following are reported:

\begin{itemize}
\tightlist
\item
  \textbf{False accept:} a case whose hidden simulator state is
  non-stable is finally accepted.
\item
  \textbf{False reject:} a case whose hidden simulator state is stable
  is finally rejected.
\item
  \textbf{Investigation count:} number of cases whose initial active
  action is \texttt{INVESTIGATE}.
\item
  \textbf{Investigation rate:} investigation count divided by case
  count.
\item
  \textbf{Decision cost:} realized sum of false-accept cost,
  false-reject cost, and investigation cost.
\item
  \textbf{Recall:} final rejects among all cases whose simulator label
  is non-stable.
\item
  \textbf{Precision:} true non-stable final rejects divided by all final
  rejects.
\end{itemize}

Because the implementation always converts an investigation into a final
accept or reject decision, an \texttt{INVESTIGATE} action is not counted
as a final detection.

\subsection{9. Results}\label{results}

\subsubsection{9.1 Primary 50-case
results}\label{primary-50-case-results}

The full primary comparison is:

\begin{longtable}[]{@{}
  >{\raggedright\arraybackslash}p{(\columnwidth - 14\tabcolsep) * \real{0.0968}}
  >{\raggedleft\arraybackslash}p{(\columnwidth - 14\tabcolsep) * \real{0.1290}}
  >{\raggedleft\arraybackslash}p{(\columnwidth - 14\tabcolsep) * \real{0.1290}}
  >{\raggedleft\arraybackslash}p{(\columnwidth - 14\tabcolsep) * \real{0.1290}}
  >{\raggedleft\arraybackslash}p{(\columnwidth - 14\tabcolsep) * \real{0.1290}}
  >{\raggedleft\arraybackslash}p{(\columnwidth - 14\tabcolsep) * \real{0.1290}}
  >{\raggedleft\arraybackslash}p{(\columnwidth - 14\tabcolsep) * \real{0.1290}}
  >{\raggedleft\arraybackslash}p{(\columnwidth - 14\tabcolsep) * \real{0.1290}}@{}}
\toprule\noalign{}
\begin{minipage}[b]{\linewidth}\raggedright
Policy
\end{minipage} & \begin{minipage}[b]{\linewidth}\raggedleft
False Accept
\end{minipage} & \begin{minipage}[b]{\linewidth}\raggedleft
False Reject
\end{minipage} & \begin{minipage}[b]{\linewidth}\raggedleft
Investigations
\end{minipage} & \begin{minipage}[b]{\linewidth}\raggedleft
Investigation Rate
\end{minipage} & \begin{minipage}[b]{\linewidth}\raggedleft
Decision Cost
\end{minipage} & \begin{minipage}[b]{\linewidth}\raggedleft
Recall
\end{minipage} & \begin{minipage}[b]{\linewidth}\raggedleft
Precision
\end{minipage} \\
\midrule\noalign{}
\endhead
\bottomrule\noalign{}
\endlastfoot
Baseline & 9 & 3 & 0 & 0.0\% & 99.0 & 0.609 & 0.824 \\
Binary & 7 & 3 & 0 & 0.0\% & 79.0 & 0.696 & 0.842 \\
\textbf{Active} & \textbf{1} & \textbf{3} & \textbf{11} &
\textbf{22.0\%} & \textbf{27.8} & \textbf{0.957} & \textbf{0.880} \\
\end{longtable}

The active-policy decision-cost calculation is visible directly:

{[} 1(10.0)+3(3.0)+11(0.8)=27.8. {]}

Relative to the static baseline in this 50-case synthetic run, active
decision cost is reduced from 99.0 to 27.8, a reduction of 71.9\%. False
accepts fall from 9 to 1, an 88.9\% reduction. Recall rises from 0.609
to 0.957, while the active policy performs 11 investigations.

The binary belief policy also improves on the static threshold, reducing
false accepts from 9 to 7 and decision cost from 99.0 to 79.0 without
any investigations. This isolates a useful point: probabilistic
reasoning alone changes the result even before active information
collection is added.

\subsubsection{9.2 Confusion-style counts for the primary 50-case
run}\label{confusion-style-counts-for-the-primary-50-case-run}

The final action confusion counts are printed below.
\texttt{drift\_event=True} means the simulator state is not
\texttt{STABLE}.

\paragraph{Baseline}\label{baseline}

\begin{longtable}[]{@{}lrr@{}}
\toprule\noalign{}
Hidden evaluation label & Final ACCEPT & Final REJECT \\
\midrule\noalign{}
\endhead
\bottomrule\noalign{}
\endlastfoot
Stable & 24 & 3 \\
Non-stable & 9 & 14 \\
\end{longtable}

\paragraph{Binary}\label{binary}

\begin{longtable}[]{@{}lrr@{}}
\toprule\noalign{}
Hidden evaluation label & Final ACCEPT & Final REJECT \\
\midrule\noalign{}
\endhead
\bottomrule\noalign{}
\endlastfoot
Stable & 24 & 3 \\
Non-stable & 7 & 16 \\
\end{longtable}

\paragraph{Active}\label{active}

\begin{longtable}[]{@{}lrr@{}}
\toprule\noalign{}
Hidden evaluation label & Final ACCEPT & Final REJECT \\
\midrule\noalign{}
\endhead
\bottomrule\noalign{}
\endlastfoot
Stable & 24 & 3 \\
Non-stable & 1 & 22 \\
\end{longtable}

The 50-case run therefore contains 23 non-stable and 27 stable cases.

\subsubsection{9.3 120-case before/after
re-test}\label{case-beforeafter-re-test}

The retained pre-change and post-change results are:

\begin{longtable}[]{@{}
  >{\raggedright\arraybackslash}p{(\columnwidth - 20\tabcolsep) * \real{0.0698}}
  >{\raggedleft\arraybackslash}p{(\columnwidth - 20\tabcolsep) * \real{0.0930}}
  >{\raggedleft\arraybackslash}p{(\columnwidth - 20\tabcolsep) * \real{0.0930}}
  >{\raggedleft\arraybackslash}p{(\columnwidth - 20\tabcolsep) * \real{0.0930}}
  >{\raggedleft\arraybackslash}p{(\columnwidth - 20\tabcolsep) * \real{0.0930}}
  >{\raggedleft\arraybackslash}p{(\columnwidth - 20\tabcolsep) * \real{0.0930}}
  >{\raggedleft\arraybackslash}p{(\columnwidth - 20\tabcolsep) * \real{0.0930}}
  >{\raggedleft\arraybackslash}p{(\columnwidth - 20\tabcolsep) * \real{0.0930}}
  >{\raggedleft\arraybackslash}p{(\columnwidth - 20\tabcolsep) * \real{0.0930}}
  >{\raggedleft\arraybackslash}p{(\columnwidth - 20\tabcolsep) * \real{0.0930}}
  >{\raggedleft\arraybackslash}p{(\columnwidth - 20\tabcolsep) * \real{0.0930}}@{}}
\toprule\noalign{}
\begin{minipage}[b]{\linewidth}\raggedright
Policy
\end{minipage} & \begin{minipage}[b]{\linewidth}\raggedleft
Before Cost
\end{minipage} & \begin{minipage}[b]{\linewidth}\raggedleft
After Cost
\end{minipage} & \begin{minipage}[b]{\linewidth}\raggedleft
Before False Accept
\end{minipage} & \begin{minipage}[b]{\linewidth}\raggedleft
After False Accept
\end{minipage} & \begin{minipage}[b]{\linewidth}\raggedleft
Before Investigations
\end{minipage} & \begin{minipage}[b]{\linewidth}\raggedleft
After Investigations
\end{minipage} & \begin{minipage}[b]{\linewidth}\raggedleft
Before Recall
\end{minipage} & \begin{minipage}[b]{\linewidth}\raggedleft
After Recall
\end{minipage} & \begin{minipage}[b]{\linewidth}\raggedleft
Before Precision
\end{minipage} & \begin{minipage}[b]{\linewidth}\raggedleft
After Precision
\end{minipage} \\
\midrule\noalign{}
\endhead
\bottomrule\noalign{}
\endlastfoot
Baseline & 242.0 & 242.0 & 23 & 23 & 0 & 0 & 0.549 & 0.549 & 0.875 &
0.875 \\
Binary & 222.0 & 222.0 & 21 & 21 & 0 & 0 & 0.588 & 0.588 & 0.882 &
0.882 \\
\textbf{Active} & \textbf{222.0} & \textbf{136.8} & \textbf{21} &
\textbf{10} & \textbf{0} & \textbf{31} & \textbf{0.588} & \textbf{0.804}
& \textbf{0.882} & \textbf{0.911} \\
\end{longtable}

The post-change active policy therefore:

\begin{itemize}
\tightlist
\item
  reduces decision cost from 222.0 to 136.8, a reduction of 38.4\%;
\item
  reduces false accepts from 21 to 10, a reduction of 52.4\%;
\item
  raises recall from 0.588 to 0.804, an increase of 0.216 or 21.6
  percentage points;
\item
  raises precision from 0.882 to 0.911, an increase of 0.029 or 2.9
  percentage points;
\item
  increases investigation count from 0 to 31, moving the investigation
  rate from 0.0\% to 25.8\%.
\end{itemize}

There are 51 non-stable and 69 stable cases in the 120-case evaluation.

\subsubsection{9.4 Post-change 120-case confusion
counts}\label{post-change-120-case-confusion-counts}

\begin{longtable}[]{@{}lrr@{}}
\toprule\noalign{}
Hidden evaluation label & Final ACCEPT & Final REJECT \\
\midrule\noalign{}
\endhead
\bottomrule\noalign{}
\endlastfoot
Stable & 65 & 4 \\
Non-stable & 10 & 41 \\
\end{longtable}

The pre-change active policy had 21 final false accepts and 4 final
false rejects, while the post-change active policy has 10 false accepts
and 4 false rejects. The improvement therefore comes primarily from
converting accepted non-stable cases into rejected cases rather than
reducing false rejection.

\subsubsection{9.5 Investigation-cost
sensitivity}\label{investigation-cost-sensitivity}

The primary 50-case benchmark was rerun with investigation costs of 0.4,
0.6, 0.8, 1.0, and 1.2.

\begin{longtable}[]{@{}
  >{\raggedleft\arraybackslash}p{(\columnwidth - 14\tabcolsep) * \real{0.1250}}
  >{\raggedleft\arraybackslash}p{(\columnwidth - 14\tabcolsep) * \real{0.1250}}
  >{\raggedleft\arraybackslash}p{(\columnwidth - 14\tabcolsep) * \real{0.1250}}
  >{\raggedleft\arraybackslash}p{(\columnwidth - 14\tabcolsep) * \real{0.1250}}
  >{\raggedleft\arraybackslash}p{(\columnwidth - 14\tabcolsep) * \real{0.1250}}
  >{\raggedleft\arraybackslash}p{(\columnwidth - 14\tabcolsep) * \real{0.1250}}
  >{\raggedleft\arraybackslash}p{(\columnwidth - 14\tabcolsep) * \real{0.1250}}
  >{\raggedleft\arraybackslash}p{(\columnwidth - 14\tabcolsep) * \real{0.1250}}@{}}
\toprule\noalign{}
\begin{minipage}[b]{\linewidth}\raggedleft
Investigation cost
\end{minipage} & \begin{minipage}[b]{\linewidth}\raggedleft
Investigations
\end{minipage} & \begin{minipage}[b]{\linewidth}\raggedleft
Investigation rate
\end{minipage} & \begin{minipage}[b]{\linewidth}\raggedleft
False Accept
\end{minipage} & \begin{minipage}[b]{\linewidth}\raggedleft
False Reject
\end{minipage} & \begin{minipage}[b]{\linewidth}\raggedleft
Decision Cost
\end{minipage} & \begin{minipage}[b]{\linewidth}\raggedleft
Recall
\end{minipage} & \begin{minipage}[b]{\linewidth}\raggedleft
Precision
\end{minipage} \\
\midrule\noalign{}
\endhead
\bottomrule\noalign{}
\endlastfoot
0.4 & 11 & 22.0\% & 1 & 3 & 23.4 & 0.957 & 0.880 \\
0.6 & 11 & 22.0\% & 1 & 3 & 25.6 & 0.957 & 0.880 \\
0.8 & 11 & 22.0\% & 1 & 3 & 27.8 & 0.957 & 0.880 \\
1.0 & 10 & 20.0\% & 1 & 3 & 29.0 & 0.957 & 0.880 \\
1.2 & 10 & 20.0\% & 1 & 3 & 31.0 & 0.957 & 0.880 \\
\end{longtable}

Within this tested range, the qualitative detection metrics remain
unchanged while investigation count changes slightly and realized cost
rises with the investigation price. This supports the use of 0.8 as a
documented synthetic cost assumption, but it does not establish that 0.8
is an optimal real-world cost.

\subsection{10. Failure Analysis}\label{failure-analysis}

\subsubsection{10.1 Five failures from the 120-case
re-test}\label{five-failures-from-the-120-case-re-test}

The project was required to classify at least five failures from the new
run. The five representative cases are:

\begin{longtable}[]{@{}
  >{\raggedright\arraybackslash}p{(\columnwidth - 12\tabcolsep) * \real{0.1364}}
  >{\raggedright\arraybackslash}p{(\columnwidth - 12\tabcolsep) * \real{0.1364}}
  >{\raggedright\arraybackslash}p{(\columnwidth - 12\tabcolsep) * \real{0.1364}}
  >{\raggedright\arraybackslash}p{(\columnwidth - 12\tabcolsep) * \real{0.1364}}
  >{\raggedright\arraybackslash}p{(\columnwidth - 12\tabcolsep) * \real{0.1364}}
  >{\raggedleft\arraybackslash}p{(\columnwidth - 12\tabcolsep) * \real{0.1818}}
  >{\raggedright\arraybackslash}p{(\columnwidth - 12\tabcolsep) * \real{0.1364}}@{}}
\toprule\noalign{}
\begin{minipage}[b]{\linewidth}\raggedright
Case
\end{minipage} & \begin{minipage}[b]{\linewidth}\raggedright
State
\end{minipage} & \begin{minipage}[b]{\linewidth}\raggedright
Failure class
\end{minipage} & \begin{minipage}[b]{\linewidth}\raggedright
Initial action
\end{minipage} & \begin{minipage}[b]{\linewidth}\raggedright
Final action
\end{minipage} & \begin{minipage}[b]{\linewidth}\raggedleft
Cost
\end{minipage} & \begin{minipage}[b]{\linewidth}\raggedright
Evidence of failure
\end{minipage} \\
\midrule\noalign{}
\endhead
\bottomrule\noalign{}
\endlastfoot
\texttt{case-066} & \texttt{DEGRADED} & Aggregate-quality blind spot &
\texttt{ACCEPT} & \texttt{ACCEPT} & 10.0 & Quality 0.842, semantic drift
0.075, latency 379.9 ms, no extra probe. \\
\texttt{case-087} & \texttt{DISTRIBUTION\_SHIFT} & Distribution-shift
under-detection & \texttt{ACCEPT} & \texttt{ACCEPT} & 10.0 & Quality
0.807 and semantic drift 0.460 produce a posterior with
\texttt{DISTRIBUTION\_SHIFT} 0.260, but the direct \texttt{DEGRADED}
probability is only 0.136. \\
\texttt{case-039} & \texttt{DISTRIBUTION\_SHIFT} & Probe non-resolution
& \texttt{INVESTIGATE} & \texttt{ACCEPT} & 10.8 & Capability audit
fires, but the post-probe posterior still leads to \texttt{ACCEPT}. \\
\texttt{case-027} & \texttt{STABLE} & Quality-threshold false alarm &
\texttt{REJECT} & \texttt{REJECT} & 3.0 & Quality 0.672 pushes the
degraded posterior to 0.412 despite low semantic drift of 0.017. \\
\texttt{case-080} & \texttt{STABLE} & Measurement-noise false alarm &
\texttt{INVESTIGATE} & \texttt{REJECT} & 3.8 & Quality 0.612 and
semantic drift 0.466 create strong degraded-looking evidence; the
capability probe does not reverse the decision. \\
\end{longtable}

\subsubsection{10.2 Failure 1: aggregate-quality blind
spot}\label{failure-1-aggregate-quality-blind-spot}

\texttt{case-066} is the clearest residual blind spot. It is truly
\texttt{DEGRADED}, yet the measured quality is 0.842 and semantic drift
is only 0.075. Its posterior after the available evidence is:

\begin{longtable}[]{@{}lr@{}}
\toprule\noalign{}
State & Posterior \\
\midrule\noalign{}
\endhead
\bottomrule\noalign{}
\endlastfoot
\texttt{STABLE} & 0.857 \\
\texttt{DEGRADED} & 0.034 \\
\texttt{TRANSIENT} & 0.041 \\
\texttt{DISTRIBUTION\_SHIFT} & 0.067 \\
\end{longtable}

The active policy therefore accepts it. This failure cannot be fixed by
merely lowering the decision threshold without creating more stable
false alarms. The missing ingredient is a capability-specific
observation.

\subsubsection{10.3 Failure 2: distribution-shift
under-detection}\label{failure-2-distribution-shift-under-detection}

\texttt{case-087} is genuinely \texttt{DISTRIBUTION\_SHIFT}. Its
posterior is:

\begin{longtable}[]{@{}lr@{}}
\toprule\noalign{}
State & Posterior \\
\midrule\noalign{}
\endhead
\bottomrule\noalign{}
\endlastfoot
\texttt{STABLE} & 0.472 \\
\texttt{DEGRADED} & 0.136 \\
\texttt{TRANSIENT} & 0.132 \\
\texttt{DISTRIBUTION\_SHIFT} & 0.260 \\
\end{longtable}

The most likely state is \texttt{STABLE}, and \texttt{P(DEGRADED)}
remains below 0.230769, so the direct cost rule accepts the case. This
identifies a policy-model mismatch: the evaluation labels all non-stable
states as drift events, but the direct cost model only prices
\texttt{DEGRADED} as the dangerous hidden state.

\subsubsection{10.4 Failure 3: probe
non-resolution}\label{failure-3-probe-non-resolution}

\texttt{case-039} is a scheduled capability audit. The probe is
purchased, but the final outcome is still \texttt{ACCEPT}. The final
posterior is:

\begin{longtable}[]{@{}lr@{}}
\toprule\noalign{}
State & Posterior \\
\midrule\noalign{}
\endhead
\bottomrule\noalign{}
\endlastfoot
\texttt{STABLE} & 0.194 \\
\texttt{DEGRADED} & 0.183 \\
\texttt{TRANSIENT} & 0.207 \\
\texttt{DISTRIBUTION\_SHIFT} & 0.415 \\
\end{longtable}

This is an important negative result. Investigation is not equivalent to
successful detection. A probe has value only if its outcome changes the
downstream decision often enough to justify its cost.

\subsubsection{10.5 Failure 4: quality-threshold false
alarm}\label{failure-4-quality-threshold-false-alarm}

\texttt{case-027} is \texttt{STABLE}, but quality falls to 0.672. The
posterior becomes:

\begin{longtable}[]{@{}lr@{}}
\toprule\noalign{}
State & Posterior \\
\midrule\noalign{}
\endhead
\bottomrule\noalign{}
\endlastfoot
\texttt{STABLE} & 0.254 \\
\texttt{DEGRADED} & 0.412 \\
\texttt{TRANSIENT} & 0.211 \\
\texttt{DISTRIBUTION\_SHIFT} & 0.123 \\
\end{longtable}

The policy rejects because 0.412 is above 0.230769. The model interprets
one noisy quality signal as strong evidence of persistent degradation.
This motivates explicit evaluator-noise modeling and calibration.

\subsubsection{10.6 Failure 5: measurement-noise false
alarm}\label{failure-5-measurement-noise-false-alarm}

\texttt{case-080} is \texttt{STABLE}, but the initial evidence gives a
very high \texttt{DEGRADED} posterior. After the capability audit, the
final posterior remains:

\begin{longtable}[]{@{}lr@{}}
\toprule\noalign{}
State & Posterior \\
\midrule\noalign{}
\endhead
\bottomrule\noalign{}
\endlastfoot
\texttt{STABLE} & 0.0014 \\
\texttt{DEGRADED} & 0.5959 \\
\texttt{TRANSIENT} & 0.3481 \\
\texttt{DISTRIBUTION\_SHIFT} & 0.0546 \\
\end{longtable}

The final action is therefore \texttt{REJECT}, producing a false-reject
cost of 3.0 plus the 0.8 investigation cost, for a total of 3.8. The
probe does not repair a strongly misleading initial observation.

\subsubsection{10.7 What the failures
imply}\label{what-the-failures-imply}

The failures separate into three broad mechanisms:

\begin{enumerate}
\def\labelenumi{\arabic{enumi}.}
\tightlist
\item
  \textbf{Missing evidence:} the monitor cannot detect a capability
  failure that is invisible to aggregate quality and generic runtime
  checks.
\item
  \textbf{Loss-model mismatch:} distribution shift can be a real change
  without being treated as a directly harmful \texttt{DEGRADED} state by
  the current utility function.
\item
  \textbf{Evaluator uncertainty:} a poor observation can be produced by
  a stable case, so the probability model needs noise calibration and
  possibly repeated measurements.
\end{enumerate}

The redesign therefore addresses one specific gap but does not claim to
have solved black-box monitoring in general.

\subsection{2. Introduction}\label{introduction}

Hosted LLM services create a monitoring problem that differs from
conventional software observability. A conventional service can often
expose version identifiers, internal state, structured error codes, and
deterministic components. A hosted language model may instead expose a
stable API name while its behavior changes over time. Chen, Zaharia, and
Zou documented substantial behavioral variation between March and June
2023 versions of GPT-3.5 and GPT-4 across tasks including mathematics,
coding, instruction following, and safety-sensitive questions. Their
results show why one-time evaluation is not enough when the serving
system is externally observable but internally opaque. {[}1{]}

The practical problem addressed here is narrower than ``detect model
drift'' in general. It is the following decision problem:

\begin{quote}
When the true state of a black-box LLM is hidden and evidence is
incomplete, can a monitor use probabilistic belief and bounded
additional evidence to reduce costly false acceptance compared with a
static quality threshold?
\end{quote}

The project is designed as an executable prototype rather than as a
production monitoring claim. The latent states, priors, likelihoods,
costs, and probe schedule are explicitly exposed so that each assumption
can be challenged.

The central engineering choice is to separate three questions:

\begin{enumerate}
\def\labelenumi{\arabic{enumi}.}
\tightlist
\item
  \textbf{What should the monitor believe?} Use a probabilistic belief
  over hidden states.
\item
  \textbf{What information should it seek next?} Estimate the expected
  value of an additional probe.
\item
  \textbf{When should it stop collecting evidence?} Compare information
  value against the probe cost and impose a hard two-round budget.
\end{enumerate}

The experiment then asks whether those choices change measurable
outcomes on a controlled synthetic testbed.

\subsection{3. Related Work}\label{related-work}

\subsubsection{3.1 LLM behavior drift}\label{llm-behavior-drift}

Chen, Zaharia, and Zou introduced LLMDrift to study how the behavior of
hosted models changes over time. Their repeated evaluations showed that
the same model service can vary across task families and that behavior
changes are not limited to a single aggregate quality score. {[}1{]}
This project shares the black-box monitoring premise but moves from
measurement toward action selection: it asks not only whether a signal
moved, but whether the monitor should accept, investigate, or reject.

The distinction between input distribution shift and model behavior
change is also operationally important. A change in inputs does not
imply a provider-side model update, while a provider-side update can
change outputs even when the input distribution is stable. The
four-state simulator therefore keeps \texttt{DEGRADED} and
\texttt{DISTRIBUTION\_SHIFT} as separate hypotheses even though the
current loss model does not yet price them symmetrically.

\subsubsection{3.2 Partial observability and active decision
making}\label{partial-observability-and-active-decision-making}

Partially Observable Markov Decision Processes provide a formal
framework for decision making when the true state is hidden and
observations are noisy. Lauri, Hsu, and Pajarinen survey this setting
and emphasize belief-state reasoning as a natural representation for
hidden-state control problems. {[}2{]}

This project uses only a small static approximation of that idea. It
maintains a belief over four hidden states, but it does not implement a
full POMDP solver, transition model, dynamic programming solution, or
learned policy. The active action is a one-step probe choice followed by
a direct cost-sensitive decision.

\subsubsection{3.3 LLM-as-a-judge and evaluator
reliability}\label{llm-as-a-judge-and-evaluator-reliability}

The monitoring problem depends heavily on the quality of the observation
layer. Recent work shows that automated LLM evaluators are not
automatically reliable. Fu and Liu found substantial inconsistency in
multilingual LLM-as-a-Judge settings, including an average Fleiss' kappa
of approximately 0.3 across their evaluated conditions. {[}3{]} Lee et
al.~proposed checklist-style evaluation to reduce evaluator variance and
improve interpretability. {[}4{]} Work on Agent-as-a-Judge similarly
argues that evaluating agent behavior requires more than inspecting a
final scalar outcome and can benefit from intermediate signals. {[}5{]}

These results matter directly to this project. A monitor that treats its
evaluator score as ground truth can turn evaluator noise into false
alarms or false confidence. The failure cases in Section 10 show both
directions: a stable case can look degraded, while a degraded case can
look healthy.

\subsubsection{3.4 Position of the present
prototype}\label{position-of-the-present-prototype}

The contribution is therefore not a claim of a new general-purpose
monitoring algorithm. It is an integrated, inspectable prototype that
combines:

\begin{itemize}
\tightlist
\item
  an explicit hidden-state belief model;
\item
  an asymmetric cost-derived accept/reject boundary;
\item
  one-step value-of-information probe selection;
\item
  a bounded investigation budget;
\item
  a concrete evidence-coverage intervention based on observed failures;
  and
\item
  a before/after re-test showing whether the intervention changes
  measured behavior.
\end{itemize}

The remaining validation burden is empirical calibration, repeated-seed
testing, real trace replay, evaluator agreement measurement, and
ablation of the active components.

\subsection{11. Discussion, Reproducibility, and Design
Implications}\label{discussion-reproducibility-and-design-implications}

\subsubsection{11.1 What the primary experiment
supports}\label{what-the-primary-experiment-supports}

The primary 50-case run supports three limited conclusions.

First, probabilistic reasoning changes the decision surface. The binary
belief policy reduces false accepts from 9 to 7 and decision cost from
99.0 to 79.0 compared with the static quality threshold.

Second, active evidence collection can materially change the result in
the constructed testbed. The active policy reaches 1 false accept, 3
false rejects, 11 investigations, and 27.8 total decision cost.

Third, the active intervention is valuable partly because it adds
evidence that the original observation layer did not contain. The
capability-coverage audit is therefore better understood as an
observation-layer redesign than as a mere threshold tweak.

\subsubsection{11.2 What the re-test
supports}\label{what-the-re-test-supports}

The 120-case re-test provides the clearest engineering evidence in the
project because it retains the pre-change result and executes the
post-change policy on the same seeded test setup.

The active policy moves from:

\begin{longtable}[]{@{}lrr@{}}
\toprule\noalign{}
Metric & Before & After \\
\midrule\noalign{}
\endhead
\bottomrule\noalign{}
\endlastfoot
False accepts & 21 & 10 \\
False rejects & 4 & 4 \\
Investigations & 0 & 31 \\
Investigation rate & 0.0\% & 25.8\% \\
Decision cost & 222.0 & 136.8 \\
Recall & 0.588 & 0.804 \\
Precision & 0.882 & 0.911 \\
\end{longtable}

The observed movement is consistent with the intended mechanism: more
cases receive additional evidence, and fewer non-stable cases are
accepted.

However, the re-test does \textbf{not} isolate the causal contribution
of each code change. The investigation-accounting correction and the
capability-coverage audit were introduced together. The largest change
in behavior is expected to come from the new audit, but the experiment
does not separately estimate the effect of the threshold correction.

\subsubsection{11.3 Why cost should be treated as a scenario
parameter}\label{why-cost-should-be-treated-as-a-scenario-parameter}

The sensitivity run shows that changing the investigation cost from 0.4
to 1.2 changes realized decision cost from 23.4 to 31.0 and changes
investigations from 11 to 10, while recall and precision remain 0.957
and 0.880 across all five tested values.

This is useful for the model because it means the architecture does not
depend on one fragile scalar. At the same time, the cost values remain
hypothetical. A real deployment should estimate investigation cost from
API charges, latency, reviewer time, opportunity cost, and the business
impact of delaying a final action.

\subsubsection{11.4 Reproducibility
procedure}\label{reproducibility-procedure}

The final project was re-executed with:

\begin{Shaded}
\begin{Highlighting}[]
\NormalTok{pytest {-}q}
\end{Highlighting}
\end{Shaded}

The result was:

\begin{Shaded}
\begin{Highlighting}[]
\NormalTok{6 passed in 0.05s}
\end{Highlighting}
\end{Shaded}

The primary benchmark was executed with:

\begin{Shaded}
\begin{Highlighting}[]
\NormalTok{python {-}m experiments.run\_experiment}
\end{Highlighting}
\end{Shaded}

using 50 cases and seed 7.

The 120-case re-test was executed with:

\begin{Shaded}
\begin{Highlighting}[]
\NormalTok{python {-}m experiments.run\_scaled}
\end{Highlighting}
\end{Shaded}

using 120 cases and seed 7.

The investigation-cost sensitivity experiment was executed with:

\begin{Shaded}
\begin{Highlighting}[]
\NormalTok{python {-}m experiments.sensitivity}
\end{Highlighting}
\end{Shaded}

using investigation costs 0.4, 0.6, 0.8, 1.0, and 1.2.

The project therefore contains executable evidence for the reported
tables rather than only manually typed results.

\subsection{12. Limitations and Knowledge
Gaps}\label{limitations-and-knowledge-gaps}

The most important unresolved questions are listed explicitly because
each one can change the interpretation of the results.

\begin{longtable}[]{@{}
  >{\raggedright\arraybackslash}p{(\columnwidth - 6\tabcolsep) * \real{0.2500}}
  >{\raggedright\arraybackslash}p{(\columnwidth - 6\tabcolsep) * \real{0.2500}}
  >{\raggedright\arraybackslash}p{(\columnwidth - 6\tabcolsep) * \real{0.2500}}
  >{\raggedright\arraybackslash}p{(\columnwidth - 6\tabcolsep) * \real{0.2500}}@{}}
\toprule\noalign{}
\begin{minipage}[b]{\linewidth}\raggedright
Knowledge gap
\end{minipage} & \begin{minipage}[b]{\linewidth}\raggedright
Why it matters
\end{minipage} & \begin{minipage}[b]{\linewidth}\raggedright
Current status
\end{minipage} & \begin{minipage}[b]{\linewidth}\raggedright
Required next test
\end{minipage} \\
\midrule\noalign{}
\endhead
\bottomrule\noalign{}
\endlastfoot
Prior calibration & Posterior quality depends on the starting belief. &
Prior 0.72/0.12/0.08/0.08 is hypothetical. & Estimate priors from
labeled incident history or replay traces. \\
Likelihood calibration & Synthetic likelihoods can make the Bayesian
model look stronger or weaker than it really is. & All likelihoods are
hand-specified. & Fit likelihoods from labeled observations and measure
calibration. \\
Conditional dependence & Correlated evidence can lead naive Bayes to
overcount evidence. & Independence is assumed, not tested. & Compare
naive Bayes against a calibrated discriminative or graphical model. \\
Distribution-shift validation & The simulator contains shift states, but
real shift signatures are not validated. & Synthetic only. & Add
input-distribution features and real replay labels. \\
Multi-state loss function & Current direct decision cost is based only
on \texttt{P(DEGRADED)}. & Distribution shift can be detected
probabilistically but is not directly priced. & Define state-specific
action losses and compare them with a scalar degraded-risk rule. \\
Probe quality & Investigation is useful only if the probe is selective
enough to change decisions. & Five failure classes include probe
non-resolution. & Measure probe sensitivity, specificity, and
decision-change rate. \\
Probe cost and latency & 0.8 is a synthetic action cost, not a measured
operational cost. & Source-specific time and monetary cost are missing.
& Log API cost, wall-clock latency, compute, and human-review time for
each probe. \\
Audit frequency & Every fourth observation is a hand-designed schedule.
& Not learned or optimized. & Compare intervals such as 2, 4, 8 and
adaptive schedules under repeated seeds. \\
VOI attribution & The current re-test changes VOI accounting and adds
the audit together. & Causal contribution is confounded. & Run binary,
random-probe, VOI-probe, and audit-only ablations. \\
Temporal persistence & The current belief is reinitialized per case. &
No temporal filtering across cases. & Add a state-transition model and
persistent belief updates. \\
Repeated-seed uncertainty & One seed cannot quantify variance across
synthetic workloads. & Primary and scaled runs use seed 7. & Repeat
across many seeds and report confidence intervals. \\
Evaluator reliability & A noisy evaluator can create false alarms or
miss failures. & Known failure classes demonstrate this risk. & Compare
automated evaluation against human or reference labels and measure
agreement. \\
Historical replay & Synthetic labels are controlled but artificial. &
Historical replay path is not yet part of the experiment. & Replay real
incidents with frozen prompts and model settings where available. \\
Human escalation & The code does not emit a human-review action. & Human
review is treated as an operational boundary. & Add explicit review cost
and a three-way final policy if the deployment requires it. \\
Probe adversarial robustness & A targeted probe can itself be sensitive
to framing, prompt, or context. & Not tested. & Evaluate probe stability
under paraphrase and adversarial inputs. \\
\end{longtable}

\subsubsection{12.1 Statistical
limitation}\label{statistical-limitation}

The reported reductions are descriptive results on one deterministic
synthetic workload. The paper does not provide confidence intervals,
hypothesis tests, or a claim of generalization beyond that workload.

\subsubsection{12.2 Ground-truth
limitation}\label{ground-truth-limitation}

The simulator knows its own hidden state, which is useful for controlled
evaluation but unlike production monitoring where the true state is
usually delayed, partial, or disputed. In particular, the project does
not prove that a real \texttt{DISTRIBUTION\_SHIFT} label would be
available at decision time.

\subsubsection{12.3 Cost-model limitation}\label{cost-model-limitation}

The false-accept cost of 10.0, false-reject cost of 3.0, and
investigation cost of 0.8 define a synthetic utility surface. Different
organizations would rationally use different numbers. Changing those
values can change the decision boundary and the frequency of
investigation without changing the underlying belief model.

\subsubsection{12.4 Calibration
limitation}\label{calibration-limitation}

A posterior of 0.80 is currently a model output, not a demonstrated
statement that roughly 80\% of comparable cases are degraded.
Calibration must be measured on independent labeled data before
posterior values are interpreted operationally.

\subsection{13. Conclusion}\label{conclusion}

This project implements a compact black-box LLM monitoring agent that
treats the monitored system as a hidden-state process and combines
probabilistic belief, information theory, active evidence selection, and
cost-sensitive action.

The main engineering result is not a claim of production superiority. It
is a reproducible demonstration that changing the \textbf{information
available to the monitor} can matter more than simply moving a final
decision threshold.

In the primary 50-case synthetic benchmark, the static quality baseline
records 9 false accepts and 3 false rejects at a total decision cost of
99.0. The binary belief policy reduces those values to 7 false accepts,
3 false rejects, and 79.0 cost. The active policy records 1 false
accept, 3 false rejects, 11 investigations, and 27.8 cost.

More importantly, the 120-case before/after re-test shows a measurable
movement after the design intervention. Active decision cost falls from
222.0 to 136.8, false accepts fall from 21 to 10, investigations rise
from 0 to 31, recall increases from 0.588 to 0.804, and precision
increases from 0.882 to 0.911.

The failure analysis prevents overinterpretation. The redesigned monitor
still misses distribution shifts, can investigate without resolving
uncertainty, and can reject stable cases when the evaluator evidence is
misleading. These are not side notes. They identify the next research
tasks: empirical likelihood calibration, probe-quality measurement,
multi-state loss functions, repeated-seed evaluation, historical
incident replay, evaluator-agreement studies, and ablation of the active
components.

The appropriate conclusion is therefore bounded: \textbf{the prototype
supports the hypothesis that bounded active evidence collection can
improve cost-sensitive monitoring decisions in a controlled synthetic
black-box LLM testbed, but the evidence is not sufficient to establish
calibrated real-world performance.}

\subsection{References}\label{references}

{[}1{]} Lingjiao Chen, Matei Zaharia, and James Zou. ``How is ChatGPT's
behavior changing over time?'' arXiv:2307.09009, 2023. The study
compares March and June 2023 versions of GPT-3.5 and GPT-4 across
multiple task families and reports substantial behavior changes.
https://arxiv.org/abs/2307.09009

{[}2{]} Mikko Lauri, David Hsu, and Joni Pajarinen. ``Partially
Observable Markov Decision Processes in Robotics: A Survey.''
arXiv:2209.10342, 2022. https://arxiv.org/abs/2209.10342

{[}3{]} Xiyan Fu and Wei Liu. ``How Reliable is Multilingual
LLM-as-a-Judge?'' Findings of the Association for Computational
Linguistics: EMNLP 2025, pages 11040-11053.
https://aclanthology.org/2025.findings-emnlp.587/

{[}4{]} Yukyung Lee, JoongHoon Kim, Jaehee Kim, Hyowon Cho, Jaewook
Kang, Pilsung Kang, and Najoung Kim. ``CheckEval: A reliable
LLM-as-a-Judge framework for evaluating text generation using
checklists.'' EMNLP 2025, pages 15771-15798.
https://aclanthology.org/2025.emnlp-main.796/

{[}5{]} Mingchen Zhuge et al.~``Agent-as-a-Judge: Evaluate Agents with
Agents.'' ICML 2025. https://icml.cc/virtual/2025/poster/45485

{[}6{]} Claude E. Shannon. ``A Mathematical Theory of Communication.''
Bell System Technical Journal, 27(3), 379-423 and 27(4), 623-656, 1948.

\subsection{22. AI-use statement}\label{ai-use-statement}

OpenAI ChatGPT was used as a research and engineering assistant during
the development and revision of this project. The model was used for
specific, reviewable tasks rather than being treated as an authority on
the experimental result.

\subsubsection{Tasks supported by
ChatGPT}\label{tasks-supported-by-chatgpt}

ChatGPT was used to:

\begin{enumerate}
\def\labelenumi{\arabic{enumi}.}
\tightlist
\item
  clarify the problem formulation as a hidden-state, partially
  observable monitoring problem;
\item
  derive the cost-sensitive accept/reject threshold from the stated
  false-accept and false-reject costs;
\item
  check the Bayesian update and entropy calculations, including the
  correction that the worked evidence example increases entropy from
  1.291 bits to 1.842 bits, a change of +0.551 bits;
\item
  inspect the investigation implementation and identify the mismatch
  between gross value of information, net value of information, and the
  reported expected investigation cost;
\item
  revise the investigation rule to require
  \texttt{gross\_VOI\ \textgreater{}\ investigation\_cost} and to report
  post-probe expected decision cost plus investigation cost;
\item
  review failure cases and help define the five failure classes used in
  the re-test;
\item
  propose and implement the bounded capability-coverage audit that was
  introduced because aggregate quality could remain healthy while a
  capability degraded;
\item
  reorganize and proofread the paper, including the tables and numerical
  narrative;
\item
  identify literature relevant to black-box LLM drift, partially
  observable decision making, and evaluator reliability; and
\item
  check the project for remaining methodological gaps that require
  future experiments rather than unsupported claims.
\end{enumerate}

\subsubsection{Verification performed outside the model's prose
generation}\label{verification-performed-outside-the-models-prose-generation}

The following numerical and implementation details were checked
explicitly against the project files and experiment outputs:

\begin{itemize}
\tightlist
\item
  the prior values 0.72, 0.12, 0.08, and 0.08 sum to 1.00;
\item
  the worked posterior values 21.6\%, 32.3\%, 8.7\%, and 37.3\% sum to
  100.0\% after normalization;
\item
  the entropy calculation uses 1.291 bits before evidence and 1.842 bits
  after evidence, giving +0.551 bits;
\item
  the cost-derived decision threshold is 3/13 = 0.230769;
\item
  the active 50-case cost is 1(10.0) + 3(3.0) + 11(0.8) = 27.8;
\item
  the 120-case active before/after decision costs are 222.0 and 136.8;
\item
  the 120-case false accepts are 21 before the redesign and 10 after it;
\item
  the five representative failure classifications correspond to
  \texttt{case-066}, \texttt{case-087}, \texttt{case-039},
  \texttt{case-027}, and \texttt{case-080}; and
\item
  the final automated test suite reports 6 passed tests.
\end{itemize}

\subsubsection{Experiments executed for the final
revision}\label{experiments-executed-for-the-final-revision}

The final project environment was re-executed with the following
commands:

\begin{Shaded}
\begin{Highlighting}[]
\NormalTok{pytest {-}q}
\end{Highlighting}
\end{Shaded}

Result: \textbf{6 passed}.

\begin{Shaded}
\begin{Highlighting}[]
\NormalTok{python {-}m experiments.run\_experiment}
\end{Highlighting}
\end{Shaded}

Configuration: \textbf{50 cases, seed 7}, comparing baseline, binary,
and active policies.

\begin{Shaded}
\begin{Highlighting}[]
\NormalTok{python {-}m experiments.run\_scaled}
\end{Highlighting}
\end{Shaded}

Configuration: \textbf{120 cases, seed 7}, retaining the pre-change
active results for before/after comparison.

\begin{Shaded}
\begin{Highlighting}[]
\NormalTok{python {-}m experiments.sensitivity}
\end{Highlighting}
\end{Shaded}

Configuration: investigation costs \textbf{0.4, 0.6, 0.8, 1.0, and 1.2}
on the 50-case benchmark.

\subsubsection{Human responsibility}\label{human-responsibility}

The mathematical assumptions, synthetic likelihoods, costs, simulator
design, interpretation of failures, and conclusions remain subject to
human review. ChatGPT output was treated as editable analysis and
code-review input. It was not treated as experimental ground truth. The
reported results come from the executable project and the listed
experiment runs, while the paper distinguishes measured synthetic
results from hypothetical assumptions and unresolved research questions.

\subsection{Abstract}\label{abstract}

Hosted language-model APIs can change observable behavior while exposing
little information about the provider-side cause of that change. This
paper presents a small active monitor for that setting. The monitor
maintains a probabilistic belief over four hidden states,
\texttt{STABLE}, \texttt{DEGRADED}, \texttt{TRANSIENT}, and
\texttt{DISTRIBUTION\_SHIFT}, and chooses among \texttt{ACCEPT},
\texttt{INVESTIGATE}, and \texttt{REJECT} under asymmetric synthetic
costs. A naive-Bayes-style evidence model combines task quality,
semantic drift, formatting, safety, latency, and API-error signals. A
one-step value-of-information rule determines whether a discretionary
probe is worth its fixed investigation cost, while a hard two-round
budget prevents unbounded evidence collection.

The main engineering correction was to separate gross value of
information from investigation cost and to report the actual expected
post-probe decision cost. Failure analysis then exposed a structural
observation gap: some degraded cases maintained healthy aggregate
quality and runtime signals. A bounded capability-coverage audit was
introduced every fourth observation within each task category.

On the deterministic 50-case benchmark, the static quality baseline
produced 9 false accepts, 3 false rejects, and total decision cost 99.0.
The binary belief policy produced 7 false accepts, 3 false rejects, and
cost 79.0. The redesigned active policy produced 1 false accept, 3 false
rejects, 11 investigations, and cost 27.8. In the 120-case before/after
re-test, active decision cost fell from 222.0 to 136.8 and false accepts
fell from 21 to 10, while recall increased from 0.588 to 0.804.

The results are synthetic and use hypothetical priors, likelihoods, and
costs. The paper therefore does not claim calibrated probabilities or
production superiority. Its main result is narrower: in a controlled
black-box monitoring testbed, bounded additional evidence can change
cost-sensitive decisions materially, and failure analysis shows that
better decision rules cannot compensate for evidence that the
observation layer does not contain.

\end{document}
